\documentclass[11pt]{article}

\usepackage[margin=1in]{geometry}
\usepackage{amsmath, amssymb}
\usepackage{booktabs}
\usepackage[hidelinks]{hyperref}
\usepackage[htt]{hyphenat}
\sloppy

\newcommand{\Afd}{A_{\mathrm{FD}}}
\newcommand{\Afe}{A_{\mathrm{FE}}}
\newcommand{\R}{\mathbb{R}}
\DeclareMathOperator{\cond}{\kappa}

\title{The PINO loss as a least-squares problem:\\ structure, conditioning, preconditioning}
\date{\today}

\begin{document}
\maketitle

\noindent
Working notes, not paper material. The question: PINO's \texttt{mse} reduction looks like it should
be \emph{formally} the same object as our bare least-squares arm, differing only in which discrete
Laplacian appears. This confirms that, identifies the matrix precisely, works out its spectrum and
conditioning, asks how one would precondition it, and then treats the \texttt{rel} reduction that
PINO actually ships --- whose geometry is not least-squares at all.

Every number below is computed at the paper's resolution, $n_x=n_y=65$ on the unit square
($h = 1/64$, $N=63^2=3969$ interior nodes), with the repo's own assembly
(\texttt{PoissonProblem}) for the finite-element operator. Analytic and numerical spectra agree to
six digits, so the closed forms quoted here can be trusted.

\section{Setup}

The PDE is $-\Delta u = f$ on $\Omega=(0,1)^2$ with $u|_{\partial\Omega}=0$. Write $u\in\R^N$ for
the vector of interior nodal values and $f\in\R^N$ for the source. Both arms impose the boundary
condition \emph{hard, in the model}: our arms through \texttt{apply\_zero\_boundary}, PINO (since
the switch away from the mollifier) through \texttt{ZeroBoundaryModel}, which is the same
operation. So in both cases the boundary values are exactly zero and can be eliminated, leaving a
square system on the interior.

\paragraph{Our least-squares arm.} With $\Afe$ the $Q_1$ stiffness matrix and $M$ the consistent
mass matrix,
\begin{equation}
  L_{\mathrm{LS}}(u) \;=\; \tfrac12 \bigl\| \Afe u - M f \bigr\|_2^2 .
  \label{eq:ls}
\end{equation}

\paragraph{PINO, \texttt{mse} reduction.} The Laplacian is taken with second-order central
differences and the residual is averaged over interior points only:
\begin{equation}
  L_{\mathrm{mse}}(u)
  \;=\; \frac{1}{N}\sum_{i,j \,\mathrm{interior}} \bigl( -\Delta_h u - f \bigr)_{ij}^2
  \;=\; \frac{1}{N}\bigl\| \Afd u - f \bigr\|_2^2 ,
  \label{eq:mse}
\end{equation}
where $\Afd$ is the five-point Laplacian
$(\Afd u)_{ij} = h^{-2}\,(4u_{ij}-u_{i\pm1,j}-u_{i,j\pm1})$.

\section{The \texttt{mse} variant is a least-squares problem}

Comparing \eqref{eq:ls} and \eqref{eq:mse}: \textbf{yes, formally identical in structure} ---
both are $\tfrac{c}{2}\|Au-b\|_2^2$ for a symmetric positive-definite $A$ and a fixed right-hand
side, on a single sample. The differences are entirely in \emph{which} $(A,b,c)$.

That statement can be sharpened, because the five-point stencil is itself a finite-element
operator. Scaling out the $h^{-2}$,
\begin{equation}
  h^2 \Afd \;=\;
  \begin{bmatrix} & -1 & \\ -1 & 4 & -1 \\ & -1 & \end{bmatrix}
  \;=\; A_{P_1},
\end{equation}
the stiffness matrix of \emph{piecewise-linear $P_1$ elements on the uniform right-triangulation}
of the same grid (verified numerically: $h^2\Afd$ has entries $4$ and $-1$ exactly). Likewise,
$f$ sampled pointwise is $h^{-2}$ times the \emph{lumped} load $M_L f$ with $M_L = h^2 I$ ---
the consistent mass matrix $M$ has interior row sums exactly $h^2$, so lumping replaces $Mf$ by
$h^2 f$. Therefore
\begin{equation}
  \boxed{\;
  L_{\mathrm{mse}}(u) \;=\; \frac{h^{-4}}{N}\,\bigl\| A_{P_1} u - M_L f \bigr\|_2^2 \;}
  \label{eq:identification}
\end{equation}
and PINO's \texttt{mse} loss is exactly our least-squares loss with three substitutions:
\begin{enumerate}
  \item $Q_1$ bilinear quadrilaterals $\to$ $P_1$ linear triangles (nine-point $\to$ five-point);
  \item the consistent mass matrix $\to$ a lumped one (equivalently: exact quadrature of the load
        $\to$ pointwise sampling);
  \item an overall constant $\tfrac12 \to h^{-4}/N$.
\end{enumerate}
At $n=65$ that constant differs by a factor $2h^{-4}/N \approx 8.5\times10^{3}$. Under plain
gradient descent this would demand a correspondingly smaller step; under Adam it very largely
cancels, since Adam divides by a running RMS of the gradient. This is worth remembering when
reading the learning-rate sweeps: the reductions differ by orders of magnitude in scale and the
sweeps still land within a factor of a few of one another, precisely because the optimizer is
scale-free.

The third substitution is bookkeeping. The first two are real, and they are what the rest of this
note is about.

\section{Properties of \texorpdfstring{$\Afd$}{A\_FD}}

\subsection{Structure}

$\Afd$ is symmetric positive definite, block-tridiagonal with tridiagonal blocks, and has the
Kronecker form $\Afd = h^{-2}(T\otimes I + I\otimes T)$ with $T=\mathrm{tridiag}(-1,2,-1)$. It is
block-Toeplitz-Toeplitz-block: the stencil is constant across the grid, which is a consequence of
the uniform mesh and would be lost on a graded or unstructured one. Bandwidth is $n_x-2$;
five nonzeros per interior row against nine for $\Afe$.

\subsection{Spectrum: both operators share an eigenbasis}

On a uniform grid with homogeneous Dirichlet conditions the discrete sine functions
\begin{equation}
  v^{(k,l)}_{ij} \;=\; \sin(k\pi i h)\,\sin(l\pi j h), \qquad k,l = 1,\dots,n_x-2,
\end{equation}
are \emph{exact} eigenvectors of both operators. This is not an approximation and not special to
finite differences: it follows from the constant stencil plus the odd reflection at the boundary.
Verified numerically --- applying either matrix and comparing against the diagonal action in
DST-I space gives a relative error of $5.9\times10^{-16}$ ($\Afe$) and $3.4\times10^{-16}$
($\Afd$). Consequently both are diagonalised by the same fast sine transform, and every spectral
statement below is available in closed form. With $\theta_k = k\pi h$:
\begin{align}
  \lambda^{\mathrm{FD}}_{kl}
    &= \frac{4}{h^2}\Bigl( \sin^2\tfrac{\theta_k}{2} + \sin^2\tfrac{\theta_l}{2} \Bigr), \\
  \lambda^{\mathrm{FE}}_{kl}
    &= \tfrac13\Bigl( 8 - 2\cos\theta_k - 2\cos\theta_l - 4\cos\theta_k\cos\theta_l \Bigr).
\end{align}

\subsection{Conditioning}

\begin{center}
\begin{tabular}{lrrr}
\toprule
& $\lambda_{\min}$ & $\lambda_{\max}$ & $\cond(A)$ \\
\midrule
$\Afe$ ($Q_1$, 9-point) & $4.816\times10^{-3}$ & $3.997$              & $829.9$  \\
$\Afd$ (5-point)        & $19.74$              & $3.275\times10^{4}$  & $1659.4$ \\
\bottomrule
\end{tabular}
\end{center}

Both scale as $\cond(A) = \Theta(h^{-2})$, as any second-order elliptic discretisation must. The
interesting part is the constant: at this resolution
\begin{equation}
  \frac{\cond(\Afd)}{\cond(\Afe)} \;=\; 1.9996 \;\approx\; 2 ,
\end{equation}
i.e.\ \textbf{the five-point Laplacian is twice as ill-conditioned as the $Q_1$ stiffness matrix}.
The reason is visible in the symbols: at the highest frequency $\theta=\pi$ the FD symbol reaches
$8h^{-2}$ while the $Q_1$ symbol saturates at $4$ (in its own scaling) --- the nine-point stencil's
corner couplings flatten the top of the spectrum, while the bottom of the spectrum is fixed for
both by consistency with $-\Delta$. Note also the enormous difference in \emph{absolute} scale:
$\Afd = \Theta(h^{-2})$ against $\Afe = \Theta(1)$, which is the $h^{-4}$ in
\eqref{eq:identification} showing up again.

For optimisation what matters is the Hessian of the loss. For a linear model
$\nabla^2 L_{\mathrm{LS}} = A^\top A = A^2$, so
\begin{equation}
  \cond(\nabla^2 L) \;=\; \cond(A)^2 \;=\;
  \begin{cases}
    6.89\times10^{5} & \text{($Q_1$ weak form)} \\
    2.75\times10^{6} & \text{(five-point FD)}
  \end{cases}
\end{equation}
a factor $4$ apart. The square is the same mechanism that motivates preconditioning in the paper;
the strong form simply makes it worse.

\section{How one would precondition \texorpdfstring{$\Afd$}{A\_FD}}

Everything we do for $\Afe$ transfers, because $\Afd$ is the more classical object of the two ---
textbook multigrid was developed for exactly this operator.

\paragraph{Geometric multigrid.} A V-cycle with damped-Jacobi smoothing, bilinear prolongation and
$R=P^\top$ applies unchanged. Only the smoother parameter differs. The optimal damping is
$\omega_{\mathrm{opt}} = 2/(\mu^{\mathrm{HF}}_{\min} + \mu_{\max})$ where $\mu$ are the eigenvalues
of $D^{-1}A$ and HF denotes the modes not representable on the coarse grid. Computed:

\begin{center}
\begin{tabular}{lrr}
\toprule
& $\omega_{\mathrm{opt}}$ & smoothing factor $\mu$ \\
\midrule
$\Afe$ ($Q_1$, 9-point) & $0.8892 \approx 8/9$ & $0.333$ \\
$\Afd$ (5-point)        & $0.8002 \approx 4/5$ & $0.599$ \\
\bottomrule
\end{tabular}
\end{center}

Both reproduce the textbook values, which is a useful check on the method we used to derive $8/9$
for $Q_1$. The second column is the more interesting one: \textbf{Jacobi is a markedly worse
smoother for the five-point stencil} --- one sweep removes a factor $0.6$ of the high-frequency
error against $1/3$ for $Q_1$. Recovering the same per-cycle error reduction therefore costs
roughly twice as many smoothing steps, since $0.6^{\,2\cdot 2} \approx 0.13$ against
$(1/3)^{2\cdot2} \approx 0.012$. The nine-point stencil's corner couplings, which cost more
arithmetic per matrix-vector product, buy a genuinely better smoother.

\paragraph{Exact spectral preconditioning.} Since the eigenvectors are known and the transform is
fast, $\Afd^{-1}$ can be applied \emph{exactly} in $O(N\log N)$ by DST: transform, divide by
$\lambda^{\mathrm{FD}}_{kl}$, transform back. On a uniform square this is strictly better than
multigrid --- it is a direct solver at near-optimal cost, and $\cond(P\Afd)=1$ by construction.
The reason we do not use it, and the reason it should not tempt anyone, is that it exists only for
constant-coefficient problems on tensor-product domains with these boundary conditions. Multigrid
survives variable coefficients, unstructured meshes and general geometry; the DST does not. If a
spectral preconditioner were used in the paper it would be measuring the wrong thing, in the same
way the $\sin(\pi x)\sin(\pi y)$ mollifier was.

\paragraph{Cheaper options.} Incomplete Cholesky, or a symmetric Gauss--Seidel sweep, both apply
and both leave the $\Theta(h^{-2})$ growth intact --- they improve the constant, not the rate. Only
a multilevel method (or the exact transform) gives $h$-independent conditioning.

\section{Finite differences against the weak form}

\begin{center}
\begin{tabular}{lll}
\toprule
& $\Afd$ (strong form) & $\Afe$ (weak form) \\
\midrule
stencil            & 5-point            & 9-point ($Q_1$) \\
scaling            & $\Theta(h^{-2})$   & $\Theta(1)$ \\
$\cond$ at $n=65$  & $1659$             & $830$ \\
smoothing factor   & $0.60$             & $0.33$ \\
right-hand side    & $f$ pointwise      & $Mf$, exact quadrature \\
eigenvectors       & \multicolumn{2}{c}{discrete sines --- shared} \\
symmetry / definiteness & \multicolumn{2}{c}{SPD --- shared} \\
consistency order  & \multicolumn{2}{c}{$O(h^2)$ --- shared} \\
\bottomrule
\end{tabular}
\end{center}

Beyond the table, three differences are conceptual rather than numerical.

\paragraph{Variational characterisation.} $\Afe u = Mf$ is the Galerkin condition: $u_h$ is the
$H^1_0$-orthogonal projection of the true solution onto the finite-element space, and the residual
$\Afe u - Mf$ is the gradient of the energy
$E(u)=\tfrac12 u^\top \Afe u - u^\top M f$. This is what makes the Deep Ritz arm possible at all,
and it is what gives the natural $B$-norm in which the preconditioned residual should be measured.
The finite-difference residual has no such characterisation --- $\Afd$ is a consistent
approximation of the operator, not the representation of a bilinear form in a basis, and there is
no energy whose gradient it is. (One may of course \emph{reverse-engineer} one via the $P_1$
identification \eqref{eq:identification}, which is exactly what that identity says; but the strong
form as written does not supply it.)

\paragraph{The right-hand side.} $Mf$ integrates the source against the basis functions; pointwise
$f$ samples it. For smooth $f$ both are second-order accurate and the difference is a lower-order
perturbation. For rough or high-frequency $f$ they are not equivalent, and pointwise sampling is
the more fragile of the two --- it can alias. With $K=10$ modes on a $65^2$ grid we are far from
that regime, so this difference should not be visible in our experiments; it would be at $K$
comparable to $n_x/2$.

\paragraph{Boundary treatment.} Because we now zero boundary nodes in both arms, the Dirichlet
elimination is identical and this is no longer a difference. Under the mollifier it was: the
mollified model could not represent any function that failed to vanish like
$\sin(\pi x)\sin(\pi y)$ near the boundary, which is a constraint on the hypothesis class, not on
the discretisation.

\section{The \texttt{rel} reduction: not a least-squares problem}

What PINO actually ships, and what we are running, is
\begin{equation}
  L_{\mathrm{rel}}(u) \;=\; \frac{\bigl\| \Afd u - f \bigr\|_2}{\|f\|_2}
  \qquad\text{(averaged over the batch),}
\end{equation}
a \emph{norm}, not a squared norm. This single square root changes the geometry completely, and the
change is not a matter of scaling.

\subsection{The loss surface is a cone}

Write $r = \Afd u - f$. Then $L_{\mathrm{rel}} = \|r\|/\|f\|$ is positively homogeneous of degree
one in $r$: along any ray towards the minimiser the loss decreases \emph{linearly}, not
quadratically. Near $r=0$ the surface is a cone with a point at the bottom, not a paraboloid, and
$L$ is not differentiable at the minimum.

The gradient in residual space is
\begin{equation}
  \frac{\partial L_{\mathrm{rel}}}{\partial r} \;=\; \frac{r}{\|r\|\,\|f\|}
  \qquad\Longrightarrow\qquad
  \Bigl\| \frac{\partial L_{\mathrm{rel}}}{\partial r} \Bigr\| \;=\; \frac{1}{\|f\|}
  \quad\text{\emph{exactly}, for every } r\neq0 .
\end{equation}
Measured across six decades of $\|r\|$ (from $6.3\times10^{1}$ down to $6.2\times10^{-5}$) the
gradient norm is constant at $1.5799\times10^{-2} = 1/\|f\|$ to all printed digits. In parameter
space,
\begin{equation}
  \frac{\partial L_{\mathrm{rel}}}{\partial u} = \frac{\Afd\, r}{\|r\|\,\|f\|},
  \qquad
  \Bigl\|\frac{\partial L_{\mathrm{rel}}}{\partial u}\Bigr\|
  = \frac{\|\Afd \hat r\|}{\|f\|}
  \in \Bigl[ \frac{\lambda_{\min}}{\|f\|}, \frac{\lambda_{\max}}{\|f\|} \Bigr],
\end{equation}
so it is bounded away from zero but direction-dependent --- it does not vanish at the solution
either.

\paragraph{Consequence.} A method with a fixed step size \emph{cannot converge} on this loss. The
gradient does not shrink as the residual does, so the iterate overshoots by a fixed distance and
settles into a limit cycle whose radius is proportional to the step size --- the same phenomenon
measured earlier for Adam on the quadratic loss ($r_\infty \approx 0.011\,\mathrm{lr}$), but here
it is forced by the loss rather than by the optimizer's variance estimate. Convergence
\emph{requires} annealing the learning rate, which the cosine schedule supplies. Anyone running
this loss at constant learning rate should expect a floor and should not read it as the method's
accuracy.

\subsection{Conditioning}

The Hessian in residual space is
\begin{equation}
  \frac{\partial^2 L_{\mathrm{rel}}}{\partial r^2}
  \;=\; \frac{1}{\|r\|\,\|f\|}\bigl( I - \hat r \hat r^\top \bigr),
  \qquad \hat r = r/\|r\| ,
\end{equation}
and pulling back to $u$,
\begin{equation}
  \nabla^2_u L_{\mathrm{rel}}
  \;=\; \frac{1}{\|r\|\,\|f\|}\, \Afd \bigl( I - \hat r \hat r^\top \bigr) \Afd .
  \label{eq:relhess}
\end{equation}
Three things follow, and they matter more than the $\cond(\Afd)^2$ of the \texttt{mse} case.

\begin{enumerate}
  \item \textbf{It is singular.} The projector kills the radial direction, so
  $\nabla^2 L_{\mathrm{rel}}$ always has a null direction. The condition number, taken literally,
  is infinite at every point. The meaningful quantity is the conditioning \emph{on the tangent
  space}.
  \item \textbf{On that tangent space the conditioning is unchanged.} Restricted to
  $\hat r^\perp$, \eqref{eq:relhess} is $\Afd^2$ up to the scalar $1/(\|r\|\|f\|)$, so the spread
  of curvatures is still $\cond(\Afd)^2 \approx 2.75\times10^{6}$. \textbf{Normalising by $\|f\|$
  does not precondition anything.} It rescales the loss and flattens one direction; the
  ill-conditioning inherited from the discrete Laplacian is untouched. Preconditioning the
  \texttt{rel} loss still requires a $P \approx \Afd^{-1}$ inside the norm, exactly as for the
  squared loss.
  \item \textbf{The scale diverges as training proceeds.} The prefactor $1/\|r\|$ grows without
  bound as $r\to0$. So the loss becomes \emph{stiffer} the closer one gets, while the gradient norm
  stays fixed --- the ratio of curvature to gradient blows up like $1/\|r\|$. This is the precise
  sense in which the \texttt{rel} loss behaves like normalised gradient descent, and it explains
  the empirical observation that the mollified PINO arm was nearly insensitive to the learning rate
  across three decades ($0.32$--$0.48$ best validation error for
  $\mathrm{lr}\in[10^{-4},3\times10^{-2}]$): the loss self-normalises the gradient scale, so the
  learning rate mostly sets the limit-cycle radius rather than the rate of progress.
\end{enumerate}

\subsection{The per-sample normalisation}

Dividing by $\|f\|$ \emph{per sample} before averaging over the batch equalises the contribution of
samples with different source amplitudes. Our summed convention does not: a sample with a large
$\|f\|$ dominates the batch gradient. Which is preferable is a modelling choice rather than a
mathematical fact --- but it is a third axis on which the two objectives differ, independent of
strong-versus-weak form and of finite differences versus assembly, and it should not be conflated
with either.

\section{Summary}

\begin{itemize}
  \item PINO's \texttt{mse} reduction \emph{is} a least-squares loss, and specifically
  \eqref{eq:identification}: our least-squares loss with $P_1$ triangles instead of $Q_1$ quads, a
  lumped instead of a consistent mass matrix, and a constant $h^{-4}/N$ in front.
  \item Its matrix $\Afd$ shares an eigenbasis with $\Afe$ (discrete sines, exactly), is SPD, and
  is $\Theta(h^{-2})$-conditioned like $\Afe$ --- but with twice the constant, hence four times the
  Hessian condition number.
  \item It is preconditioned the same way. Geometric multigrid applies verbatim with
  $\omega_{\mathrm{opt}}=4/5$ instead of $8/9$; damped Jacobi is a distinctly worse smoother on the
  five-point stencil ($0.60$ against $0.33$).
  \item The \texttt{rel} reduction that PINO actually uses is \emph{not} least squares. It is
  conically non-smooth at the minimum, its gradient norm is exactly $1/\|f\|$ everywhere, it cannot
  converge at a fixed step size, and --- the point most likely to be misread --- \textbf{its
  normalisation does not improve conditioning}. On the tangent space the curvature spread is still
  $\cond(\Afd)^2$.
\end{itemize}

\end{document}
